\chapter{Introduction}
% Github link to all the code used in this report:

\section{Interstellar's AR\&D Problem}

\section{Real-World Relevance}
% Mainly space debris removal, docking would probably be useful
\section{Approach}
% simulator + neural network + evolutionary algorithm
% Tudat chosen as simulator (discretized time and space using numerical integration), as it is open-source and has alot of functionality for orbital mechanics
% various neural nets, both fixed and variable topology
% Various evolutionary algorithms: SGA, SEA and NEAT


%  Introduction
%    Why is this work relevant? may seem just funny (which it is), but has real applications: if you can dock with space debris, it becomes alot easier to deorbit, or possibly atleast stabilize/detumble. Interstellar provides a really nice scenario to work with, and alot of people are familiar with it.
%    - Briefly explain the AR&D problem in Interstellar
%    - Briefly explain the relevance of the AR&D problem
%    - Briefly explain the approach taken in this work